\chapter{Implementation}

\section{Project Structure}

The project codebase follows a modular organization with clear separation of concerns. The directory structure is organized as follows:

\lstset{basicstyle=\ttfamily\footnotesize, numbers=none, frame=none}
\begin{lstlisting}
project_root/
  data/kaggle_dataset/          # Raw and processed Hyderabad data
  src/
    ingestion/                  # Open-Meteo API data fetcher
    models/                     # Model architecture definitions
    preprocessing/              # Data cleaning and normalization
    training/                   # Training loop and optimization
    evaluation/                 # Metrics computation and plotting
    deployment/                 # ONNX conversion and Pi tools
  scripts/
    individual_trainers/        # Per-model training scripts
  outputs/hyderabad/            # All 14 model outputs and results
    <model_name>/
      model.pth                 # PyTorch checkpoint
      metrics.json              # Performance metrics
      plots/                    # Convergence, parity, timeseries plots
      config.json               # Best hyperparameters
    all_models_comparison.csv   # Master comparison table
\end{lstlisting}

\section{Data Pipeline Implementation}

\subsection{Data Fetching}

The Open-Meteo API client is implemented in \texttt{src/ingestion/india\_aq.py}. The API request is constructed with the following parameters:

\begin{itemize}
    \item \textbf{Coordinates}: Latitude 17.3850, Longitude 78.4867 (Hyderabad)
    \item \textbf{Hourly Variables}: PM$_{2.5}$, PM$_{10}$, CO, NO$_2$, SO$_2$, O$_3$, European AQI
    \item \textbf{Time Range}: One full calendar year of historical data
    \item \textbf{Timezone}: Asia/Kolkata (UTC+5:30)
\end{itemize}

The API response is parsed and stored as a CSV file with columns: \texttt{timestamp, pm2\_5, pm10, carbon\_monoxide, nitrogen\_dioxide, sulphur\_dioxide, ozone, us\_aqi}. The CSV format ensures portability and enables easy inspection with standard tools.

\subsection{Preprocessing Pipeline}

The preprocessing pipeline is implemented as a series of transformations:

\textbf{Step 1: Load and Validate.} The raw CSV is loaded using pandas. Basic validation checks confirm that all expected columns are present and that the timestamp column is monotonically increasing.

\textbf{Step 2: Handle Missing Values.} Gaps in the time series are identified by checking for missing hourly timestamps. Missing pollutant values are imputed using linear interpolation:

\begin{equation}
    y_{\text{missing}} = y_{t_1} + \frac{y_{t_2} - y_{t_1}}{t_2 - t_1} \times (t_{\text{missing}} - t_1)
\end{equation}

\textbf{Step 3: Feature-Target Split.} The PM$_{2.5}$ column is separated as the target variable $y$, while the remaining columns (PM$_{10}$, CO, NO$_2$, SO$_2$, O$_3$, AQI) serve as features $X$.

\textbf{Step 4: Normalization.} Min-Max scaling parameters ($x_{\min}$, $x_{\max}$) are computed from the training set and applied consistently to all three splits (train, validation, test). The scaler objects are serialized using pickle for use during inference.

\textbf{Step 5: Sliding Window Construction.} A sliding window of 168 hours (7 days) creates input-output pairs. For each window position $t$, the input $X[t-168:t]$ and output $y[t:t+24]$ are paired. This yields approximately 8,500 training samples from the full dataset.

\subsection{Train/Val/Test Split}

The chronological split is implemented as:

\begin{lstlisting}[language=Python]
n_total = len(dataset)
n_train = int(0.70 * n_total)
n_val   = int(0.15 * n_total)
n_test  = n_total - n_train - n_val

train_data = dataset[:n_train]
val_data   = dataset[n_train:n_train + n_val]
test_data  = dataset[n_train + n_val:]
\end{lstlisting}

\section{Model Implementation}

\subsection{Statistical Models}

Statistical models are implemented using the \texttt{statsmodels} library:

\textbf{ARIMA} uses \texttt{ARIMA(endog=train\_y, order=(5,1,2))} from \texttt{statsmodels.tsa.arima.model}. After fitting, 24-step forecasts are generated using the \texttt{forecast(steps=24)} method. The model is refit with each new observation in a rolling manner.

\textbf{SARIMA} uses \texttt{SARIMAX(endog=train\_y, order=(2,1,1), seasonal\_order=(1,1,1,24))}. The seasonal period of 24 corresponds to the diurnal cycle.

\textbf{VAR} uses \texttt{VAR(endog=train\_data)} from \texttt{statsmodels.tsa.api}. The VAR model requires all 7 features as endogenous variables. The order $p = 2$ is selected using the Akaike Information Criterion (AIC).

\subsection{Deep Learning Models}

All deep learning models are implemented in PyTorch and share a common \texttt{BaseForecaster} interface:

\begin{lstlisting}[language=Python]
class BaseForecaster(nn.Module):
    def __init__(self, input_dim=7, hidden_dim=128,
                 output_dim=24, num_layers=2):
        super().__init__()
        # Common parameters stored, architecture
        # defined by subclass
        
    def forward(self, x):
        # x: (batch, 168, 7) -> (batch, 24)
        raise NotImplementedError
\end{lstlisting}

\textbf{LSTM Implementation.} The LSTM model is defined as:

\begin{lstlisting}[language=Python]
class LSTMForecaster(BaseForecaster):
    def __init__(self, input_dim=7, hidden_dim=128,
                 output_dim=24, num_layers=2, dropout=0.2):
        super().__init__()
        self.lstm = nn.LSTM(
            input_size=input_dim,
            hidden_size=hidden_dim,
            num_layers=num_layers,
            dropout=dropout,
            batch_first=True
        )
        self.fc = nn.Linear(hidden_dim, output_dim)
    
    def forward(self, x):
        out, (h_n, c_n) = self.lstm(x)
        # Use last hidden state
        return self.fc(out[:, -1, :])
\end{lstlisting}

\textbf{GRU Implementation.} The GRU follows an identical interface with \texttt{nn.GRU} replacing \texttt{nn.LSTM}. The key difference is the absence of a separate cell state, reducing the parameter count.

\textbf{BiLSTM Implementation.} Bidirectionality is enabled by setting \texttt{bidirectional=True} in the LSTM constructor, which doubles the effective hidden dimension. The output linear layer adapts accordingly to accept $2 \times$ \texttt{hidden\_dim} features.

\textbf{CNN-LSTM Implementation.} The convolutional front-end uses two 1D convolution layers:

\begin{lstlisting}[language=Python]
self.conv1 = nn.Conv1d(in_channels=7, out_channels=64,
                        kernel_size=3, padding=1)
self.conv2 = nn.Conv1d(in_channels=64, out_channels=128,
                        kernel_size=3, padding=1)
# After convolutions:
self.lstm = nn.LSTM(input_size=128, hidden_size=128,
                     num_layers=2, batch_first=True)
\end{lstlisting}

The input tensor is transposed from $(B, 168, 7)$ to $(B, 7, 168)$ for 1D convolution (channels-first format), then transposed back for LSTM processing.

\textbf{CNN-GRU Implementation.} Identical to CNN-LSTM with \texttt{nn.GRU} replacing the LSTM layers.

\textbf{BiLSTM with Attention Implementation.} The attention mechanism is implemented as a separate module:

\begin{lstlisting}[language=Python]
class Attention(nn.Module):
    def __init__(self, hidden_dim):
        super().__init__()
        self.W = nn.Linear(hidden_dim * 2, hidden_dim)
        self.v = nn.Linear(hidden_dim, 1)
    
    def forward(self, encoder_outputs):
        # encoder_outputs: (batch, seq_len, hidden_dim*2)
        energy = torch.tanh(self.W(encoder_outputs))
        attention = F.softmax(self.v(energy), dim=1)
        context = torch.sum(attention * encoder_outputs, dim=1)
        return context, attention
\end{lstlisting}

The context vector is passed through a final linear layer to produce the 24-hour forecast.

\textbf{Transformer Implementation.} The Transformer uses PyTorch's \texttt{nn.TransformerEncoder}:

\begin{lstlisting}[language=Python]
self.pos_encoder = PositionalEncoding(d_model=128, max_len=200)
self.encoder_layer = nn.TransformerEncoderLayer(
    d_model=128, nhead=4, dim_feedforward=512,
    dropout=0.1, batch_first=True
)
self.transformer = nn.TransformerEncoder(
    self.encoder_layer, num_layers=4
)
\end{lstlisting}

\textbf{Informer and Autoformer Implementations.} These are implemented as custom architectures following the original papers, with modifications to accept the (168, 7) input shape. The ProbSparse attention in Informer is approximated by selecting the top-$u$ queries based on the sparsity measurement. Autoformer's Auto-Correlation uses FFT for efficient period-based dependency computation.

\textbf{ST-GCN Implementation.} The graph convolution uses a learned adjacency matrix:

\begin{lstlisting}[language=Python]
# Learn node embeddings
self.node_emb = nn.Parameter(
    torch.randn(num_nodes, embed_dim)
)
# Compute adjacency via spatial attention
adj = F.softmax(
    F.relu(torch.mm(self.node_emb, self.node_emb.T)),
    dim=1
)
\end{lstlisting}

Temporal modeling uses gated 1D convolutions with kernel size 3.

\section{Training Pipeline}

\subsection{Training Loop}

All deep learning models use a standardized training loop:

\begin{lstlisting}[language=Python]
def train_model(model, train_loader, val_loader,
                epochs=200, lr=1e-3, patience=20):
    optimizer = torch.optim.Adam(model.parameters(), lr=lr)
    scheduler = ReduceLROnPlateau(optimizer, mode='min',
                                   factor=0.5, patience=10)
    criterion = nn.MSELoss()
    best_val_loss = float('inf')
    patience_counter = 0
    
    for epoch in range(epochs):
        # Training phase
        model.train()
        train_loss = 0
        for X_batch, y_batch in train_loader:
            optimizer.zero_grad()
            y_pred = model(X_batch)
            loss = criterion(y_pred, y_batch)
            loss.backward()
            optimizer.step()
            train_loss += loss.item()
        
        # Validation phase
        model.eval()
        val_loss = 0
        with torch.no_grad():
            for X_batch, y_batch in val_loader:
                y_pred = model(X_batch)
                val_loss += criterion(y_pred, y_batch).item()
        
        # Early stopping check
        if val_loss < best_val_loss:
            best_val_loss = val_loss
            patience_counter = 0
            torch.save(model.state_dict(), 'best_model.pth')
        else:
            patience_counter += 1
            if patience_counter >= patience:
                break
        
        scheduler.step(val_loss)
    
    # Restore best weights
    model.load_state_dict(torch.load('best_model.pth'))
    return model
\end{lstlisting}

\subsection{Hyperparameter Optimization with Optuna}

The Optuna integration wraps the training loop within an objective function:

\begin{lstlisting}[language=Python]
def objective(trial):
    # Suggest hyperparameters
    hidden_dim = trial.suggest_categorical(
        'hidden_dim', [64, 128, 256]
    )
    num_layers = trial.suggest_int('num_layers', 1, 3)
    dropout = trial.suggest_float('dropout', 0.1, 0.5)
    lr = trial.suggest_float('lr', 1e-4, 1e-2, log=True)
    
    # Create model with suggested params
    model = ModelClass(hidden_dim, num_layers, dropout)
    
    # Train and evaluate
    model = train_model(model, train_loader, val_loader,
                       lr=lr)
    val_rmse = evaluate_model(model, val_loader)
    
    return val_rmse  # Optuna minimizes this

study = optuna.create_study(
    direction='minimize',
    sampler=TPESampler(seed=42)
)
study.optimize(objective, n_trials=20)
\end{lstlisting}

\subsection{Best Hyperparameters Found}

Table~\ref{tab:best_hparams} presents the optimal hyperparameters found by Optuna for the top-performing models.

\begin{table}[H]
\centering
\caption{Optimal Hyperparameters for Top Models}
\label{tab:best_hparams}
\begin{tabular}{lcccccl}
\toprule
\textbf{Model} & \textbf{Hidden} & \textbf{Layers} & \textbf{Dropout} & \textbf{LR} & \textbf{Epochs} \\
\midrule
Transformer & 128 & 4 & 0.10 & 9.2e-4 & 87 \\
CNN-LSTM & 128 & 2 & 0.20 & 5.1e-4 & 123 \\
GRU & 128 & 2 & 0.15 & 3.7e-4 & 145 \\
BiLSTM & 128 & 2 & 0.20 & 4.8e-4 & 98 \\
Informer & 128 & 3 & 0.10 & 7.2e-4 & 112 \\
RNN & 128 & 2 & 0.20 & 2.9e-4 & 156 \\
CNN-GRU & 128 & 2 & 0.25 & 6.3e-4 & 108 \\
BiLSTM-Attn & 128 & 2 & 0.30 & 3.1e-4 & 135 \\
LSTM & 128 & 2 & 0.20 & 5.4e-4 & 118 \\
Autoformer & 128 & 2 & 0.10 & 8.1e-4 & 94 \\
\bottomrule
\end{tabular}
\end{table}

\subsection{Reproducibility Measures}

To ensure reproducibility, the following measures are implemented:

\begin{itemize}
    \item \textbf{Fixed Random Seed}: Python's random module, NumPy, and PyTorch are all seeded with the value 42.
    \item \textbf{Deterministic Operations}: \texttt{torch.backends.cudnn.deterministic = True} and \texttt{torch.backends.cudnn.benchmark = False} are set for GPU operations.
    \item \textbf{Data Split Consistency}: The chronological split uses fixed proportions (70/15/15) computed deterministically.
    \item \textbf{Config Serialization}: All hyperparameters and training configurations are saved to \texttt{config.json} alongside each model checkpoint.
\end{itemize}

\section{ONNX Conversion}

Trained PyTorch models are converted to ONNX format for edge deployment. The conversion process:

\begin{enumerate}
    \item \textbf{Model Preparation}: The trained model is loaded and set to evaluation mode.
    \item \textbf{Example Input}: A dummy input tensor of shape (1, 168, 7) --- representing a single 7-day window --- is created.
    \item \textbf{Tracing}: \texttt{torch.onnx.export()} traces the model computation graph:
    
    \begin{lstlisting}[language=Python]
    torch.onnx.export(
        model, dummy_input, "model.onnx",
        input_names=['input'],
        output_names=['output'],
        dynamic_axes={
            'input': {0: 'batch_size'},
            'output': {0: 'batch_size'}
        },
        opset_version=14
    )
    \end{lstlisting}
    
    \item \textbf{Validation}: The exported model is validated using ONNX's built-in checker.
    \item \textbf{Inference Verification}: Outputs from the ONNX model are compared against the original PyTorch model to verify numerical equivalence (tolerance: $10^{-5}$).
\end{enumerate}

\section{Raspberry Pi Deployment}

\subsection{Hardware Specifications}

The Raspberry Pi 4 used for deployment has the following specifications:

\begin{table}[H]
\centering
\caption{Raspberry Pi 4 Hardware Specifications}
\label{tab:pi_specs}
\begin{tabular}{ll}
\toprule
\textbf{Property} & \textbf{Value} \\
\midrule
Model & Raspberry Pi 4 Model B Rev 1.5 \\
CPU & Quad-core ARM Cortex-A72 @ 2.0 GHz (overclocked) \\
RAM & 8 GB LPDDR4 \\
GPU & VideoCore VI @ 500 MHz \\
Storage & 59 GB (32 GB used) microSD \\
OS & Debian GNU/Linux 13 (trixie) arm64 \\
Kernel & Linux 6.12.75+rpt-rpi-v8 \\
Python & 3.13.5 \\
CPU Governor & Performance (all cores at max frequency) \\
L1 Cache & 128 KiB data + 192 KiB instruction per core \\
L2 Cache & 1 MiB shared \\
\bottomrule
\end{tabular}
\end{table}

The Pi runs with the \textbf{performance} CPU governor, meaning all four Cortex-A72 cores operate at their maximum frequency of 2.0 GHz continuously. This is a factory overclock from the stock 1.5 GHz, providing 33\% additional compute throughput. The memory split allocates 948 MB to the ARM CPU and 76 MB to the GPU, optimized for headless inference workloads.

\subsection{Setup}

The Raspberry Pi 4 (8 GB RAM) deployment involves:

\begin{enumerate}
    \item \textbf{OS Installation}: Raspberry Pi OS (64-bit, Debian Bullseye) is installed on a 32 GB microSD card.
    \item \textbf{Runtime Installation}: \texttt{pip install onnxruntime} installs the ARM64-optimized ONNX Runtime.
    \item \textbf{Model Transfer}: All 14 ONNX models are transferred to the Pi via SCP.
    \item \textbf{Benchmark Script}: A Python script loads each ONNX model, performs 10 warmup inferences, then measures latency over 100 inference runs.
\end{enumerate}

\subsection{Inference Benchmark}

Each model is benchmarked with a fixed input of shape (1, 168, 7). The benchmark measures:
\begin{itemize}
    \item \textbf{Mean Latency}: Average time per inference (ms)
    \item \textbf{Standard Deviation}: Variability across runs (ms)
    \item \textbf{Memory Usage}: Peak RAM consumption during inference (MB)
\end{itemize}

Results are saved to \texttt{benchmark\_results.json} for analysis.
